\documentclass[12pt,a4paper]{report}
\usepackage[margin=1in]{geometry}
\usepackage{hyperref}
\usepackage{titlesec}
\usepackage{setspace}
\usepackage{graphicx}
\usepackage{tabularx}
\usepackage{longtable}
\usepackage{enumitem}
\usepackage{booktabs}
\usepackage{array}
\usepackage{float}
\usepackage{fancyhdr}
\usepackage{tocloft}
\usepackage{xcolor}
\usepackage{listings}
\usepackage{caption}
\usepackage{subcaption}
\usepackage{pdfpages}
\usepackage{amsmath}
\usepackage{amssymb}
\usepackage{siunitx}
\usepackage{multirow}
\usepackage{footmisc}
\usepackage{csquotes}
\usepackage{algorithm}
\usepackage{algpseudocode}

\hypersetup{
  colorlinks=true,
  linkcolor=blue,
  urlcolor=blue,
  citecolor=blue
}

% Formatting
\onehalfspacing
\setlength{\parskip}{0.6em}
\setlength{\parindent}{0pt}
\titleformat{\chapter}[hang]{\bfseries\Huge}{\thechapter}{2pc}{}
\titlespacing*{\chapter}{0pt}{-20pt}{20pt}
\titlespacing*{\section}{0pt}{1.2em}{0.4em}
\titlespacing*{\subsection}{0pt}{1.0em}{0.3em}
\titlespacing*{\subsubsection}{0pt}{0.8em}{0.2em}

% Header/Footer
\pagestyle{fancy}
\fancyhf{}
\fancyhead[LE,RO]{\thepage}
\fancyhead[LO]{\leftmark}
\fancyhead[RE]{\rightmark}
\renewcommand{\headrulewidth}{0.4pt}
\renewcommand{\footrulewidth}{0pt}
\fancypagestyle{plain}{
  \fancyhf{}
  \renewcommand{\headrulewidth}{0pt}
  \renewcommand{\footrulewidth}{0pt}
}

% Table of contents spacing
\setlength{\cftbeforetoctitleskip}{-0.5em}
\setlength{\cftaftertoctitleskip}{1em}
\setlength{\cftbeforechapskip}{1em}
\setlength{\cftbeforesecskip}{0.5em}

% Code style
\lstdefinestyle{code}{
  basicstyle=\ttfamily\small,
  breaklines=true,
  frame=single,
  backgroundcolor=\color{black!2},
  keywordstyle=\color{blue},
  commentstyle=\color{gray!70},
  stringstyle=\color{green!40!black},
  numbers=left,
  numberstyle=\tiny, numbersep=6pt
}

\begin{document}

% Title Page (a. Project Title)
\begin{titlepage}
  \centering
  \vspace*{2cm}
  {\LARGE\bfseries JOBEASE: Intelligent Career Matching and Application Automation\par}
  \vspace{1cm}
  {\Large Project Synopsis (30--35 pages)\par}
  \vspace{1.5cm}
  {\large \today\par}
  \vfill
  {\large 
  \textbf{Submitted By:}\\
  [Your Name Here]\\
  Roll Number: [XXXXX]\\[1cm]
  \textbf{Supervisor:}\\
  [Supervisor Name], Department of Computer Science\\
  [Institute/University]
  }
\end{titlepage}

\newpage
\tableofcontents
\newpage
\listoffigures
\newpage
\listoftables
\newpage

% (b) Introduction
\chapter{Introduction}
\section{Background and Problem Domain}
The modern employment ecosystem is undergoing rapid digitization, yet the candidate and employer experiences remain disjointed. Candidates rely on heterogeneous sources (job boards, company sites, referrals) and must reconcile inconsistent formats, vocabularies, and submission workflows. This fragmentation forces repetitive tasks such as tailoring resumes and cover letters, answering screening questions in varying schemas, and tracking statuses without standardized feedback loops. For employers, high application volumes combine with limited structure to create noise, while Applicant Tracking Systems (ATS) impose keyword-centric heuristics that may overlook contextual fit and transferable skills.

Against this backdrop, \textit{JobEase} proposes an integrated, AI-enabled platform that reduces friction across the end-to-end lifecycle: discovery, tailoring, application, and tracking. By combining semantic search, retrieval-augmented generation (RAG), and structured skill ontologies, the system aims to elevate signal quality for both sides of the market while preserving user privacy and control over data. The platform is designed for measurable impact: compressing time-to-application, increasing interview rates, and improving employer screening precision.

\section{Problem Statement}
\subsection{Candidate Pain Points}
\begin{itemize}[leftmargin=2em]
  \item Manual tailoring of resumes and cover letters for each role
  \item Repetitive form filling across disparate application portals
  \item Limited insight into what improves success probability
  \item Friction from inconsistent application schemas and ATS keyword filters
  \item Time-consuming process with low conversion rates
\end{itemize}

\subsection{Employer Pain Points}
\begin{itemize}[leftmargin=2em]
  \item High volume with low signal-to-noise ratios in applications
  \item Overreliance on keyword heuristics that miss contextual fit
  \item Time-consuming screening with limited structured data
  \item Difficulty in identifying candidates with transferable skills
  \item Inefficient tracking of application statuses
\end{itemize}

\section{Scope}
This synopsis presents a comprehensive view of \textit{JobEase}, including the problem framing and motivation; background study and related work; clearly enumerated objectives and success criteria; a detailed methodology with mathematical foundations for matching and generation; expected outcomes and validation strategies; tools and technologies spanning clients, services, data, and ML; applications and future scope; and references formatted in an IEEE-like style. The document emphasizes mid-to-long explanatory paragraphs to elaborate design choices and their trade-offs.

% (c) Background Study
\chapter{Background Study: Related Work and Market Landscape}
\section{Existing Solutions}
Traditional job boards (e.g., Indeed, LinkedIn Jobs, ZipRecruiter) excel at aggregation but stop short of holistic guidance. Personalization typically hinges on rudimentary filters and historical clicks, which do not capture nuanced skill transfer or career intent. Applicant Tracking Systems (Workday, Greenhouse, Lever) optimize employer pipelines but rarely assist candidates beyond status notifications, leaving tailoring and evidence construction to candidates. AI resume tools provide point solutions for formatting and keyword coverage; however, they are not deeply integrated with job ingestion, semantic matching, or cross-portal submission and tracking. Consequently, candidates must orchestrate a patchwork of tools without shared context.

\section{Gap Analysis}
\begin{enumerate}[leftmargin=2em]
  \item \textbf{Lifecycle Integration}: Absence of end-to-end automation spanning discovery, tailoring, application, and tracking leaves efficiency gains unrealized.
  \item \textbf{Semantic Alignment}: Keyword-centric approaches miss transferable skills, seniority signals, and role-contextual evidence.
  \item \textbf{Grounded Generation}: Content generation without retrieval or structured validation risks hallucinations and style drift.
  \item \textbf{Closed-Loop Learning}: Lack of unified tracking and outcome signals limits personalization and continuous improvement.
\end{enumerate}

\section{Technological Underpinnings}
Transformer-based encoders, efficient vector databases, and hybrid lexical–semantic retrieval enable robust similarity search and re-ranking. RAG frameworks combine large language models with curated, domain-specific knowledge bases to ensure factuality, consistency, and controllable style. These advances, paired with layout-aware document understanding for resumes and forms, make it feasible to deliver reliable parsing, matching, and generation under real-world constraints.

% (d) Objectives
\chapter{Objectives}
\section{Primary Objectives}
\begin{enumerate}[leftmargin=2em]
  \item Deliver personalized job discovery and ranking via hybrid retrieval and preference-aware re-ranking.
  \item Automate resume tailoring and cover letter/screening response generation with retrieval-grounded constraints.
  \item Provide application autofill across diverse portals and consolidate statuses into a unified tracker.
  \item Ensure privacy-by-design and transparent user controls, including export/delete and consent governance.
  \item Build a scalable, observable, and secure architecture that supports rapid iteration and reliability.
\end{enumerate}

\section{Success Criteria (KPIs)}
\begin{itemize}[leftmargin=2em]
  \item \textbf{Discovery}: $>\SI{30}{\percent}$ click-through on top-5 recommendations
  \item \textbf{Tailoring}: $\geq \SI{70}{\percent}$ reduction in manual edits per application
  \item \textbf{Autofill}: $\geq \SI{50}{\percent}$ of fields populated accurately on first pass
  \item \textbf{Reliability}: P95 $< \SI{500}{\milli\second}$; uptime $>\SI{99.9}{\percent}$
\end{itemize}

% (e) Proposed Methodology
\chapter{Proposed Methodology}
\section{System Overview}
\begin{itemize}[leftmargin=2em]
  \item \textbf{Client}: Cross-platform mobile and web apps (Expo/React Native) providing consistent UX and secure storage.
  \item \textbf{Gateway}: Auth, routing, rate limiting, and observability; policy enforcement and request correlation.
  \item \textbf{Services}: Parsing, matching, generation, ingestion, and tracking, decomposed for modular scaling.
  \item \textbf{Data}: Relational DB (PostgreSQL) with JSONB, vector store (pgvector/Qdrant), blob storage, and Redis cache.
\end{itemize}

\section{Data Model and Evidence Graph}
Core entities include User, CandidateProfile, JobPosting, Application, and Document. Relationships encode many-to-many links between jobs and skills, and versions for tailored documents enable auditability. An evidence graph links generated bullets and claims to underlying experiences and metrics, facilitating verifiability and editor review.

\section{NLP Pipeline}
\subsection{Preprocessing}
Tokenization, stop-word removal, lemmatization, language detection, and named entity recognition.

\subsection{Information Extraction}
Skill, experience, education, location, and salary extraction are performed with layout-aware models for robust resume parsing. Temporal normalization captures tenure and recency; section segmentation identifies achievements versus responsibilities. Confidence scores drive human-in-the-loop review for low-signal segments.

\section{Semantic Matching and Ranking}
\subsection{Embedding-based Similarity}
Let \(e_{c} \in \mathbb{R}^{d}\) denote the candidate profile embedding and \(e_{j} \in \mathbb{R}^{d}\) a job description embedding. The cosine similarity is
\[\mathrm{cos}(e_{c}, e_{j}) = \frac{e_{c}^{\top} e_{j}}{\lVert e_{c} \rVert\, \lVert e_{j} \rVert}.\]

\subsection{Hybrid Retrieval}
Combine lexical BM25 and vector similarity with interpolation parameter \(\alpha \in [0,1]\):
\[ S(\text{job}\mid\text{candidate}) = \alpha\, S_{\mathrm{BM25}} + (1-\alpha)\, \mathrm{cos}(e_{c}, e_{j}). \]

\begin{algorithm}
\caption{Hybrid Retrieval and Ranking}
\begin{algorithmic}[1]
\Procedure{HybridRank}{$candidate, jobs, \alpha$}
  \State $q_{lex} \gets BuildLexQuery(candidate.skills)$
  \State $q_{vec} \gets Embed(candidate.summary)$
  \State $C \gets BM25Search(q_{lex}, jobs)$
  \State $V \gets VectorSearch(q_{vec}, jobs)$
  \ForAll{$job$}
    \State $job.score \gets \alpha\, C[job] + (1-\alpha)\, V[job].cosine$
  \EndFor
  \State \Return $SortDescending(jobs, key=score)$
\EndProcedure
\end{algorithmic}
\end{algorithm}

\section{Retrieval-Augmented Generation (RAG)}
\begin{enumerate}[leftmargin=2em]
  \item Construct task-specific queries using candidate/job context
  \item Retrieve exemplars/templates and evidence from knowledge base
  \item Compose prompts with constraints for factual, structured outputs
  \item Generate tailored bullets/letters; validate and score
\end{enumerate}

\subsection{Constraint and Validation Layer}
Outputs are constrained for length, style, and factual alignment. Validators enforce keyword coverage, quantify impact where possible (e.g., percentages, deltas), and check reading level. Low-confidence outputs are flagged for revision, and changes are logged for traceability.

\section{Autofill and Tracking}
Schema mapping across job portals combines rule-based dictionaries and DOM-context sequence labeling to align fields. Confidence thresholds govern auto-submit vs. user-verify flows. The unified tracker consolidates statuses, reminders, and notes, enabling outcome-aware personalization and analytics.

\section{Security and Privacy}
Security follows defense-in-depth: TLS in transit, AES-256 at rest, and least-privilege access control. Privacy is embedded through data minimization, consent recording, export/delete endpoints, and PII redaction in logs. Compliance alignment with GDPR/CCPA is supported by auditable controls and incident response playbooks.

\section{Evaluation Methodology}
Offline metrics include Precision@k and nDCG for recommendations, F1 for parsing, and ROUGE/BLEU plus human ratings for generation. Online evaluations use A/B testing with guardrails (latency, error rates, safety). Ablations quantify the contribution of hybrid retrieval, prompt patterns, and knowledge base coverage.

% (f) Expected Outcomes
\chapter{Expected Outcomes}
The project will deliver a functional MVP that integrates discovery, tailoring, autofill, and tracking in a cohesive experience. From a user-impact perspective, we expect substantial reductions in time-to-application and increases in interview rates, driven by better alignment and higher-quality application materials. For employers, improved candidate signal should reduce screening costs and increase shortlist precision. The deliverables include a robust architecture with observability, documented APIs and schemas, validated ML components with both offline and online evaluations, and security and privacy controls aligned with industry expectations.

Beyond the MVP, outcome analytics and continuous learning pipelines will enable rapid iteration. Evidence-linked generation and auditing will support institutional deployments (e.g., university career services) that require accountability and review trails. Public artifacts will include technical documentation, evaluation reports, and a roadmap for employer-facing extensions.

% (g) Tools and Technologies
\chapter{Tools and Technologies}
\section{Programming Languages and Frameworks}
TypeScript across client and server enables type safety and shared contracts. React Native with Expo accelerates cross-platform delivery. Node.js/Express provides a mature ecosystem for APIs; Python may be employed for specialized ML services where libraries and tooling warrant.

\section{Datastores and Search}
PostgreSQL with JSONB balances relational integrity and flexible nested fields. Redis supports caching, rate limiting, and ephemeral sessions. Vector similarity is backed by pgvector or Qdrant. Full-text search via OpenSearch or Meilisearch complements semantic retrieval for hybrid strategies.

\section{ML/NLP}
Transformer encoders produce dense representations for matching; layout-aware parsers enhance robustness on resumes and forms; LLMs are bounded by retrieval and constraints to generate tailored content with verifiable evidence.

\section{DevOps and Observability}
Containerization enables portability; CI/CD enforces quality gates (lint, type, tests, SAST/DAST). Metrics, logs, and traces provide end-to-end observability. Dependency hygiene and image scanning reduce supply-chain risk; configuration is parameterized per environment.

% (h) Applications and Future Scope
\chapter{Applications and Future Scope}
\section{Applications}
For candidates, JobEase functions as a copilot across discovery, tailoring, and application management, providing guidance and automation while preserving control and editability. For employers, future-facing capabilities include structured shortlists, skills coverage analytics, and bias-aware recommendation tools. Institutions such as universities and workforce agencies can adopt the platform for cohort-level insights and service delivery at scale.

\section{Future Enhancements}
Future work spans deeper personalization with sequence and context models, expanded ATS/job-board integrations, built-in interview preparation and feedback loops, and skills-to-learning pathways that tie identified gaps to curated learning resources. Over time, employer-facing features can evolve toward collaborative hiring workspaces and explainable candidate ranking.

% (i) References
\chapter{References}
\noindent The references below are formatted in an IEEE-like style.
\begin{enumerate}[leftmargin=2em]
  \item J. Smith, "AI in Human Resources: Current Applications and Future Directions," \textit{Journal of AI Applications}, vol. 15, no. 2, pp. 45--62, 2023.
  \item A. Johnson and M. Brown, "Semantic Matching in Job Recommendation Systems," in \textit{Proc. WSDM}, 2022.
  \item European Union, "General Data Protection Regulation (GDPR)," Official Journal of the European Union, 2018.
  \item J. Devlin, M.-W. Chang, K. Lee, and K. Toutanova, "BERT: Pre-training of Deep Bidirectional Transformers for Language Understanding," in \textit{Proc. NAACL}, 2019.
  \item P. Lewis et al., "Retrieval-Augmented Generation for Knowledge-Intensive NLP Tasks," in \textit{Proc. NeurIPS}, 2020.
  \item California State Legislature, "California Consumer Privacy Act (CCPA)," 2018.
  \item Y. Liu et al., "RoBERTa: A Robustly Optimized BERT Pretraining Approach," arXiv:1907.11692, 2019.
  \item T. Mikolov, K. Chen, G. Corrado, and J. Dean, "Efficient Estimation of Word Representations in Vector Space," arXiv:1301.3781, 2013.
  \item A. Vaswani et al., "Attention is All You Need," in \textit{Proc. NeurIPS}, 2017.
  \item J. Pennington, R. Socher, and C. Manning, "GloVe: Global Vectors for Word Representation," in \textit{Proc. EMNLP}, 2014, pp. 1532--1543.
\end{enumerate}

\appendix
\chapter{Project PDF}
\label{app:project-pdf}
The following appendix includes the provided project document in full for reference.
\includepdf[pages=-,scale=0.95]{../project.pdf}

\end{document}


